\section{Limitations and future work}
\label{sec:limits}

\emph{Age and scale.} The system had been running for only a few weeks when we measured it, and one person maintains it. The mechanisms work; whether they are robust over years is not yet known.

\emph{No independent witness.} The defence against a malicious operator rests on public code and public logs that nobody else is yet watching. The most valuable next step costs nothing: a second party running the same workflow from a fork.

\emph{Unsigned source.} Origin authenticity is unproven. We see three paths: signatures by the Chamber, which only the Chamber can provide; transferable TLS proofs, which need a notary; and multiple independent collectors, which need only volunteers.

\emph{Unattested steps.} The deployment and the metrics pipeline are not attested, and the lake's history is not anchored. Moving the deployment into CI with a reusable workflow would extend provenance to the site itself and could raise the build level the platform claims~\citep{github-slsa3}.

\emph{Coverage.} The attested daily collection is a small sample of recent bills and votings. The complete record since 2003 lives in the lake, which is verifiable by recomputation but not attested. The system does not cover the Senate, committees' full proceedings, the budget or state legislatures.

\emph{Summaries.} Only a few bills have an explanation in simple language, none reviewed by a person.

\emph{Simple metrics.} Coincidence with an orientation is a blunt instrument. It weights a procedural vote like a constitutional amendment, cannot distinguish discipline from agreement, and is silent on votings without a recorded orientation.

\emph{Limited evaluation.} The measurements used only public data but were run by the project's author, recompute one of the three metric families for four years, and attempted no attack by the operator, such as attesting from a self-hosted runner, from a branch, or from a commit later force-pushed away. Each of these is a short experiment that someone independent should run. We also have no usage data and ran no study of whether readers understand the site, notice the proofs or interpret them correctly; a study with journalists and fact-checkers, the users we expect to benefit most, is the natural next evaluation.

\emph{Proofs cover the model, not the code.} The Lean proofs establish the properties of the constructions as described, under stated assumptions such as injective encodings. They do not show that the project's TypeScript and Rust implement those constructions, nor do they cover rounding in the score.

\subsection{Planned: taking the record to where people are}
\label{sec:future-social}

A verifiable website reaches people who visit websites. Most Brazilians meet political information elsewhere: 82\% of internet users use WhatsApp and 46\% use it for news~\citep{carro2026brazil}, and fact-checks that arrive after a falsehood has circulated miss much of its spread~\citep{reis2020debunked}. The next stage of the project, not yet built, is a WhatsApp channel and an automated pipeline that turns each day's collection into short posts and images for social networks.

Social formats strip away what makes the site checkable: a forwarded image has no manifest, and an automated political account looks, on its face, like the content farms it competes with. We therefore set four rules for the pipeline. Posts are generated in the same public workflow, after the attestation, from attested data only. Numbers, names and vote positions are placed by fixed templates, never written by a language model, which also keeps the project within the electoral court's rules on synthetic content~\citep{tse2026res23755}. Each post carries a link and a machine-readable code pointing to its source records and their proofs. And every post is itself content-addressed, listed in the day's manifest, attested and timestamped, so that the project can show that a fabricated ``Vigia do Povo'' image is not its own. What gets posted is chosen by a published rule, not by an editor, so that the neutrality earned on the data path is not spent at the last step. The risks remain: reach rewards attention, a channel can be suspended by its platform, and short formats leave little room for caveats.

\section{Conclusion}
\label{sec:conclusion}

Brazil has spent years arguing about whether its official records can be believed, while most citizens cannot name the people who vote in their name. Vigia do Povo is a small, concrete response to both conditions: the Chamber's votes in simple language and with context, with the proof attached to the number.

The technical result is that this is feasible today, for free, with tools built for another purpose. Content addressing, CI attestations and blockchain timestamps let an outsider with no authority publish legislative data that anyone can check, and an independent recomputation of the published figures from the official files found no difference. The honest boundary is the source: until the Chamber signs what it publishes, every proof starts one step too late.

The social result is a claim deliberately smaller than the technical one. The evidence says that information by itself rarely moves voters, that corrections struggle against falsehood, and that scorecards can backfire. Verifiability does not repeal any of that. What it offers is narrower and, we think, still worth having: a public record of the legislature that does not ask to be believed, that shows when it was changed, and that will still be checkable when its author is gone. In a country where the question ``who says so?'' has ended too many conversations, a record that answers ``check it yourself'' is a useful thing to build.
